\documentclass[17pt]{beamer}

\usepackage[utf8]{inputenc}
\usepackage[turkish]{babel}
\usepackage[T1]{fontenc}

\usepackage{amsmath}
\usepackage{amssymb}
\usepackage{amsthm}
\usepackage{enumerate}
\usepackage{tikz}
\usepackage{transparent}
\usepackage{xcolor}

\usetheme{Warsaw}
\usecolortheme{default}
\definecolor{nigdeyesili_acik}{RGB}{3, 150, 166}
\definecolor{nigdeyesili_koyu}{RGB}{143, 209, 217}
\setbeamercolor{structure}{fg=nigdeyesili_acik}

\author[Dr. Mehmet CANEVİ]{Arş.~Gör.~Dr.~M.~Canevi\inst{1}}

\institute{
    \inst{1}%
    Bilgisayar Mühendisliği\\
    Mühendislik Fakültesi
}
    
\date[2025] {Ders Notları, Ocak 2025}
    
\logo{\transparent{0.4}\includegraphics[width = 20mm]{logo}}
    
\title[Ders 4] {Diziler I}
\begin{document}
%%%%%%%%%%%%%%%%%%%%%%%%%%%%%%%%%%%%%%%%%%%%%%%%%%%%%%%%%%%%%%%%%%%%%%%%%%%%%%%%
\frame{\titlepage}
\begin{frame}[fragile]{İçidekiler}
    \tableofcontents
\end{frame}
%%%%%%%%%%%%%%%%%%%%%%%%%%%%%%%%%%%%%%%%%%%%%%%%%%%%%%%%%%%%%%%%%%%%%%%%%%%%%%%%
\section{Dizileri tanıyalım}
\begin{frame}[fragile]{Fibbonacci sayıları}
    Fibbonacci sayılarını tanımlayalım: 
    \begin{lstlisting}
        int fibb1=1;    
        int fibb2=1;    
        int fibb3=2;    
        int fibb4=3;    
        int fibb5=5;    
        int fibb6=8;    
        int fibb7=13;\end{lstlisting}
    Dizinin 100 terimi için 100 adet değişken gerekirdi. Bu noktada diziler kullanılır.
\end{frame}
%%%%%%%%%%%%%%%%%%%%%%%%%%%%%%%%%%%%%%%%%%%%%%%%%%%%%%%%%%%%%%%%%%%%%%%%%%%%%%%%
\begin{frame}[fragile]{Diziler}
    Diziler kullanılırsa
    \begin{lstlisting}
        int fibb[]={1,1,2,3,5,8,13};\end{lstlisting}
        ile tanımlanır. Tanımlama şekli
    \begin{lstlisting}
        veri_tipi degisken_adi[]={x1,x2,x3,x4,...};
        veri_tipi degisken_adi[3]={x1,x2,x3};
        veri_tipi degisken_adi[dizi_uzunlugu];
    \end{lstlisting}
    olarak verilebilir.
\end{frame}
%%%%%%%%%%%%%%%%%%%%%%%%%%%%%%%%%%%%%%%%%%%%%%%%%%%%%%%%%%%%%%%%%%%%%%%%%%%%%%%%
\begin{frame}[fragile]{Fibbonacci Dizileri(devam)}
    Fibbonacci dizisi 
    \begin{lstlisting}
        int fibb[5];
        int i;
        fibb[0]=1;
        fibb[1]=1;
        for(i=2;i<5;i++)
        {
            fibb[i]=fibb[i-1]+fibb[i-2];
        }
    \end{lstlisting}
\end{frame}
%%%%%%%%%%%%%%%%%%%%%%%%%%%%%%%%%%%%%%%%%%%%%%%%%%%%%%%%%%%%%%%%%%%%%%%%%%%%%%%%
\begin{frame}[fragile]{Fibbonacci Dizileri(devam)}
    Fibbonacci dizisi 
    \begin{lstlisting}
        fibb[0]=1;
        fibb[1]=1;
        fibb[2]=fibb[1]+fibb[0];
        fibb[3]=fibb[2]+fibb[1];
        fibb[4]=fibb[3]+fibb[2];
    \end{lstlisting}
\end{frame}
%%%%%%%%%%%%%%%%%%%%%%%%%%%%%%%%%%%%%%%%%%%%%%%%%%%%%%%%%%%%%%%%%%%%%%%%%%%%%%%%
\begin{frame}[fragile]{2 boyutlu diziler}
    2 boyutlu diziler
    \begin{lstlisting}
        veri_tipi degisken_adi[][]={{x11,x12,x13,...},{x21,x22,x23,...},...};
        veri_tipi degisken_adi[3][3]={{1,0,0},{0,1,0},{0,0,1}};
        veri_tipi degisken_adi[dizi_uzunlugu][dizi_uzunlugu];
    \end{lstlisting}
    ile tanımlanır.
\end{frame}
%%%%%%%%%%%%%%%%%%%%%%%%%%%%%%%%%%%%%%%%%%%%%%%%%%%%%%%%%%%%%%%%%%%%%%%%%%%%%%%%
\begin{frame}[fragile]{n boyutlu diziler}
    n boyutlu diziler
    \begin{lstlisting}
        veri_tipi degisken_adi[x][y][z]...={{{},{},...},{{},{},...},....};
        veri_tipi degisken_adi[x][y][z]...;
    \end{lstlisting}
    ile tanımlanır.
\end{frame}

%%%%%%%%%%%%%%%%%%%%%%%%%%%%%%%%%%%%%%%%%%%%%%%%%%%%%%%%%%%%%%%%%%%%%%%%%%%%%%%%
\section{Alıştırmalar}
\begin{frame}[fragile]{Sorular}
    \begin{alertblock}{Soru-1}
        Soru
    \end{alertblock}
\end{frame}
%%%%%%%%%%%%%%%%%%%%%%%%%%%%%%%%%%%%%%%%%%%%%%%%%%%%%%%%%%%%%%%%%%%%%%%%%%%%%%%%
\end{document}